\section{Contact pencils: equations, primary components, and infinitesimal layers}
\label{sec:contact-primary-structure}
Fix $d\ge4$ and $n\ge2d-1$. In this section the divisor $D$ is
fixed, although it need not be reduced. Set
$G_D^\circ=\Gr(2,E_{n,D})^\circ$. It is a smooth irreducible variety
of dimension $2d-4$. The conductor construction embeds it as the
open generating locus in the Grassmannian of subseries containing
$W_D$. Let $\mathcal J_m$ be the zeroth Fitting ideal of
$\operatorname{coker}(\Sym^m\mathcal U\to V_{mn})$ on this variety.

Trivialize $\mathcal O_D(n)$, and put $R=H^0(\mathcal O_D)$.
There is a smooth surjective frame map
\begin{equation}\label{eq:contact-frame}
 \mathcal F_D=\{(a,v):a\in R^\times,\ v\in R\setminus\C\}
 \longrightarrow G_D^\circ,\qquad (a,v)\longmapsto\langle a,av\rangle.
\end{equation}
It is the open part of the frame bundle where the first basis vector
is a unit. Its relative dimension is four. All calculations on
$\mathcal F_D$ below are therefore exact faithfully flat local
calculations for $G_D^\circ$, not transverse slices inferred from
pointwise ranks.

\begin{lemma}[The power-basis presentation]\label{lem:contact-power-basis}
On $\mathcal F_D$ the multiplication cokernel is presented by
\begin{equation}\label{eq:contact-krylov}
 K_m(v)=[1,v,v^2,\ldots,v^m]:\mathcal O^{m+1}
                              \longrightarrow R\otimes\mathcal O.
\end{equation}
For $m\ge d-1$, every later column is an $\mathcal O$-linear
combination of the first $d$ columns, and
\begin{equation}\label{eq:contact-stable-ideal}
 \mathcal J_m\mathcal O_{\mathcal F_D}=(\Delta_D(v)),\qquad
 \Delta_D(v)=\det[1,v,\ldots,v^{d-1}].
\end{equation}
For $2\le m<d-1$, $\mathcal J_m=0$.
\end{lemma}
\begin{proof}
By Theorem~\ref{thm:conductor-reduction}, it suffices to multiply
sections in the pencil $\langle a,av\rangle$ on $D$. Its image
is generated by $a^m v^j$ for $0\le j\le m$. Multiplication by
$a^{-m}$ is an invertible target operation and gives
\eqref{eq:contact-krylov}. Multiplication by $v$ is an endomorphism
of the rank-$d$ free module $R\otimes\mathcal O$. Its monic
characteristic polynomial annihilates it over the coefficient ring,
by the adjugate proof of Cayley--Hamilton. Thus $v^d$ and every
higher power lie in the span of the first $d$ powers, with regular
coefficients. Column operations reduce the presentation to
$[K_{d-1}(v),0]$. This proves \eqref{eq:contact-stable-ideal}, as an
identity of ideals and not merely of radicals. If $m+1<d$, all
$d$-minors are zero. Faithful flatness descends this assertion.
\end{proof}

Write $D=\sum_{i=1}^s d_i p_i$, with distinct $p_i$. Choose local
parameters $z_i$ and the Chinese-remainder identification
\[
 R=\prod_{i=1}^s\C[z_i]/(z_i^{d_i}),\qquad
 v_i=\lambda_i+\sum_{r=1}^{d_i-1}c_{ir}z_i^r.
\]
All $\lambda_i,c_{ir}$ are independent affine coordinates before
removing the scalar locus. For $d_i\ge2$ define $T_i$ by failure
of $A$ to separate $2p_i$; for $i<j$ define $E_{ij}$ by failure
to separate $p_i+p_j$. On the frame space they have equations
$c_{i1}=0$ and $\lambda_j-\lambda_i=0$, respectively.

\begin{theorem}[Full fixed-contact classification]\label{thm:contact-factorization}
For every $m\ge d-1$, in the ordered jet basis $1,z_i,\ldots,z_i^{d_i-1}$,
\begin{equation}\label{eq:contact-confluent-determinant}
 \Delta_D(v)=
 \prod_{i:d_i\ge2}c_{i1}^{\binom{d_i}{2}}
 \prod_{i<j}(\lambda_j-\lambda_i)^{d_i d_j}.
\end{equation}
The complete primary decomposition on $G_D^\circ$ is
\begin{equation}\label{eq:contact-full-primary}
 \mathcal J_m=
 \bigcap_{i:d_i\ge2}\mathcal I_{T_i}^{\binom{d_i}{2}}
 \cap\bigcap_{i<j}\mathcal I_{E_{ij}}^{d_i d_j}.
\end{equation}
All displayed divisors are nonempty, smooth, irreducible, and
distinct. They are precisely the associated subvarieties; no
embedded associated points occur. The failure divisor has dimension
$2d-5$ and divisor class $\binom d2$ times the restricted Pluecker
class.

Every multiplication corank, for every $m\ge2$, is determined as
follows. Put $e_i=\operatorname{ord}_{z_i}(v_i-\lambda_i)$, with
$e_i=d_i$ for a constant $v_i$, and put $\ell_i=\lceil d_i/e_i\rceil$.
For each distinct value $\lambda$ among the $\lambda_i$, set
$L_\lambda=\max_{i:\lambda_i=\lambda}\ell_i$, and let
$s(v)=\sum_\lambda L_\lambda$. Then
\begin{equation}\label{eq:contact-all-coranks}
 \operatorname{corank}(\Sym^m U_A\longrightarrow V_{mn})
                 =d-\min\{m+1,s(v)\}.
\end{equation}
On the frame space every higher Fitting ideal is, exactly, the
corresponding minor ideal of \eqref{eq:contact-krylov}.
\end{theorem}
\begin{proof}
For a polynomial $P$ of degree at most $d-1$, Taylor expansion gives
\[
 P(v_i(z_i))=\sum_{r=0}^{d_i-1}
             \frac{P^{(r)}(\lambda_i)}{r!}
                    (v_i(z_i)-\lambda_i)^r\pmod {z_i^{d_i}}.
\]
The linear change from the divided derivatives of $P$ at $\lambda_i$
to the coefficients of $P(v_i(z_i))$ is triangular, with diagonal
$1,c_{i1},\ldots,c_{i1}^{d_i-1}$. Its determinant is
$c_{i1}^{\binom{d_i}{2}}$, including when this coefficient is zero.
The determinant of the divided-derivative evaluation map
\[
 P\longmapsto\left(\frac{P^{(r)}(\lambda_i)}{r!}\right)_{i,\,0\le r<d_i}
\]
in the basis $1,T,\ldots,T^{d-1}$ is the confluent Vandermonde
$\prod_{i<j}(\lambda_j-\lambda_i)^{d_i d_j}$. Here is an explicit
normalization of that identity. Start with $d$ distinct variables,
cluster $d_i$ of them near $\lambda_i$, and in each cluster replace
successive evaluation rows by divided-difference rows. The ordinary
Vandermonde determinant is divided by the product of the differences
\emph{within} each cluster. In the cluster limit the rows are the
divided derivatives above, while the remaining cross-cluster
factors tend to the displayed product. This proves the identity on
a dense open set and hence as a polynomial identity, including all
collisions. Combining the two determinants proves
\eqref{eq:contact-confluent-determinant} without inverting any of
its factors.

For any length-two $Z\le D$, restriction $E_{n,D}\to E_{n,Z}$ is
surjective. The locus in $\Gr(2,E_{n,D})$ where its restriction to
$A$ has rank at most one is the corresponding irreducible Schubert
divisor. Its rank-zero locus is excluded by generation on $D$.
Its rank-one locus is smooth: choose a basis of $A$ with one nonzero
image, and vary the second basis vector in a complementary target
direction; the determinant then has a nonzero linear derivative.
For irreducibility one may also mark the image line in $E_{n,Z}$;
the resulting incidence is a Grassmannian bundle over $\Pj(E_{n,Z})$,
and projects birationally onto the rank-one locus. The generating
open is nonempty for each divisor when $d\ge4$. On frames this
follows directly by choosing $a=1$ and a nonscalar $v$ with the
single prescribed linear equality. The same coordinates show that
the divisors are distinct.

On each affine frame chart the coordinate ring is a localization of
a polynomial ring. The distinct linear factors in
\eqref{eq:contact-confluent-determinant} generate distinct prime
ideals there. Unique factorization gives
$(\prod f_\alpha^{b_\alpha})=\bigcap(f_\alpha^{b_\alpha})$;
each ideal $(f_\alpha^{b_\alpha})$ is primary. These identities
descend by faithful flatness to \eqref{eq:contact-full-primary}.
Equivalently, at each point of the smooth base the local ring is
regular, each displayed divisor is smooth, and its powers are
primary. This also proves absence of embedded primes. The dimension
follows since this is a nonzero Cartier divisor. Each component is
a Pluecker hyperplane section, and
\[
 \sum_i\binom{d_i}{2}+\sum_{i<j}d_i d_j=\binom d2,
\]
which proves the class formula.

Finally, in $\C[z_i]/(z_i^{d_i})$ the element $v_i-\lambda_i$
has nilpotency index $\ell_i$. Its minimal polynomial is
$(T-\lambda_i)^{\ell_i}$. In the product algebra the minimal
polynomial is therefore
\begin{equation}\label{eq:contact-minimal-polynomial}
 P_v(T)=\prod_\lambda(T-\lambda)^{L_\lambda}.
\end{equation}
No polynomial of smaller degree can annihilate $v$, and division by
this monic polynomial shows that $1,v,\ldots,v^{s(v)-1}$ form a
basis of $\C[v]$. Consequently the rank of $K_m(v)$ is
$\min\{m+1,s(v)\}$. The cokernel isomorphism of
Theorem~\ref{thm:conductor-reduction} gives
\eqref{eq:contact-all-coranks} and the assertion about Fitting ideals.
\end{proof}

\begin{theorem}[Exact infinitesimal layers and singularities]
\label{thm:contact-layers}
Let $m\ge d-1$. At a point of a smooth frame chart discard the
nonvanishing factors of \eqref{eq:contact-confluent-determinant}.
Write the remaining distinct linear factors as $f_\alpha$, with
multiplicities $b_\alpha$, and put
\[
 F=\prod_\alpha f_\alpha^{b_\alpha},\qquad P=\prod_\alpha f_\alpha.
\]
In the regular local ring $B$ of that chart the failure ring is
$B/(F)$ and its nilradical is $\mathcal N=(P)/(F)$. Its exact
nilpotency index is $\max_\alpha b_\alpha$, with index one for a
reduced local ring. For every $q\ge1$, there is an isomorphism
of $B$-modules, with indicated generator $P^q$,
\begin{equation}\label{eq:contact-nilradical-layers}
 \mathcal N^q/\mathcal N^{q+1}
 \simeq
 B\big/\left(P,\prod_\alpha f_\alpha^{(b_\alpha-q)_+}\right).
\end{equation}
The completed local equation is the same weighted linear
arrangement, with unused frame and higher-jet coordinates as smooth
parameters. Its tangent cone is defined by the homogeneous product
of the active linear factors with these same multiplicities.
Its singular scheme on this chart has ideal
$(F,\partial F)$, and its singular support is the union of the
components with $b_\alpha\ge2$ and the pairwise intersections of
distinct components. These descriptions descend smoothly to
$G_D^\circ$.
\end{theorem}
\begin{proof}
The local ring in question is a localization of a polynomial ring
in the unit coordinates of $a$ and the independent jet coordinates
of $v$. The factors $f_\alpha$ are prime, distinct, and vanish at
the selected point. Unique factorization shows that $P^q\in(F)$
if and only if $q\ge b_\alpha$ for every $\alpha$. This proves the
nilpotency assertion, including at multiple intersections.

The homomorphism $B\to\mathcal N^q/\mathcal N^{q+1}$ sending
$1$ to the class of $P^q$ is onto. Its kernel is
\[
 ((P^{q+1},F):P^q)
       =(P)+(F:P^q)
       =\left(P,\prod_\alpha f_\alpha^{(b_\alpha-q)_+}\right).
\]
For the middle equality, if $rP^q=sP^{q+1}+tF$, subtract $sP$
from $r$; cancellation gives membership in $(F:P^q)$. The last
colon is computed by prime-factor exponents. This proves
\eqref{eq:contact-nilradical-layers}. Globally the chosen generator
carries the corresponding line-bundle trivialization; the displayed
isomorphism is a local one and makes no claim of a preferred global
generator.

At the chosen point every active factor is linear in centered
coordinates. Every discarded factor is a unit, so the completed
ideal is exactly the displayed product, not a truncation of it.
The initial ideal is generated by the same homogeneous product,
which proves the tangent-cone statement. The hypersurface
cotangent presentation gives $(F,\partial F)$. At a point on
exactly one component of multiplicity one the differential is a
unit times $df_\alpha\ne0$. At every point on a multiple component
or on two distinct components it is zero. This proves the singular
support assertion. Smooth faithfully flat pullback preserves these
local statements; in particular the four frame parameters do not
alter the nilpotency index or the singular-support criterion.
\end{proof}

For illustration, length four already gives subseries of codimension
two, in every $V_n$ with $n\ge7$. For every $m\ge3$ their fixed-contact
primary equations are as follows. Put $\delta_{ij}=\lambda_j-\lambda_i$;
only first-jet coefficients of multiple points are displayed.
\begin{center}
\begin{tabular}{@{}llc@{}}
\toprule
Multiplicities of $D$ & Exact equation on frames & Largest local nilpotency index\\
\midrule
$(1,1,1,1)$ & $\prod_{i<j}\delta_{ij}$ & $1$\\
$(2,1,1)$ & $c_{11}\delta_{12}^2\delta_{13}^2\delta_{23}$ & $2$\\
$(2,2)$ & $c_{11}c_{21}\delta_{12}^4$ & $4$\\
$(3,1)$ & $c_{11}^3\delta_{12}^3$ & $3$\\
$(4)$ & $c_{11}^6$ & $6$\\
\bottomrule
\end{tabular}
\end{center}
The reduced arrangement need not have normal crossings: three
coincident values give three braid hyperplanes in a two-dimensional
normal space. The equation, rather than a normal-crossing label,
records that singularity. Formula~\eqref{eq:contact-all-coranks}
also distinguishes points with the same support of failure but
different higher-jet vanishing orders.
